\section{Introduction}

滑らかな射影曲面を Lagrangian torus の族として記述する問題では，
正則部分の action coordinates と，退化に現れる divisor の組合せ論を
大域的に結び付ける必要がある．
本稿の中心結果は，一般位置の全平面退化において，actual 負 Kähler--Einstein 形式に関する
大域 Lagrangian fibration と，その射影に標識付き基底収縮を合成した
同じ写像によって，径の尺度での計量極限を実現することである．
この構成を支える相対幾何として，正の重み付き SNC divisor を持つ射影曲面の有限相対 compactification を構成し，
それを一般型超曲面の退化へ適用する．
曲線成分には任意の genus を許し，divisor 上の marked Morse 葉層，
core 近傍の primitive，交点の product germs を同時に制御する．
この相対構成から，次数 $d\geq5$ の滑らかな超曲面
$X_d\subset\PP^3$ に対し，hyperplane class の正の倍数 $aH$ に属する
任意の Kähler form の下で，全正則 fibers が Lagrangian $T^2$ となる fibration を得る．
さらに，$d=e_1+\cdots+e_k$ という次数分割を選ぶことで，core の個数と型，
node 数，および底空間の homotopy type を変更できる．
五次・六次の全平面分割では，focus affine singularities を残した大域作用適合 matroidal atlas まで構成する．
また混合五次分割では，個別の core charts の存在と，任意の branch denominators を許した大域接続の不可能性を示す．
全平面五次で任意の奇数分母を許す局所 chart の整数構造と，
共有 branch に沿う大域的な算術制約は，companion 論文
\cite{QuinticArithmetic}で扱う．

固定した一般位置の全平面退化に対し，actual 負 Kähler--Einstein 形式で
同じ相対構成を計量的に制御する．$R_t=\log(-\log|t|^2)$ とすると，
標識付き基底収縮と fibration の合成 $f_t=p_T\pi_t$ が
\[
 \sup_{x,y\in X_t}
 \left|R_t^{-1}d_{g_t}(x,y)-d_\square(f_t(x),f_t(y))\right|\to0
\]
を満たすことを，任意の固定 $d\ge5$ で証明する
（定理\ref{ke:global-realization}，定理\ref{ke:global-metric-realization}）．
これは全曲面の ambient 距離についての主張であり，nodal，branch，core fibers と，
六次以上で基底収縮が潰す内部木上の全 fiber 合併を含む．
一般の symplectic 存在構成と，この特定退化族に対する KE 計量実現は，
仮定と証明を分けて記述する．関連する局所 KE 崩壊および
Lagrangian 構成との比較は第\ref{intro:ke-collapse}節で行う．

\paragraph{三つの存在主張の量化．}\label{intro:existence-scopes}
本稿で用いる選択の自由度は，次の三つで異なる．
\begin{enumerate}
\item \emph{指定 Kähler 形式での一般存在．}
滑らかな $X_d$ と $[\omega]=aH$ の Kähler 形式を先に固定し，
実際の punctured curves 上の所定の型の marked Morse 葉層の相対 isotopy class
（全平面分割では marked pants trees）を先に指定して
fibration を構成する（定理\ref{main:global}，定理\ref{part:global}）．
SNC 相対定理では，正の weighted divisor 表示と各成分の二つ以上の交点を仮定し，
同じ Kähler class の代表 $\omega'$ と normal holonomy を構成中に選ぶ
（定理\ref{snc:main}）．
\item \emph{Actual KE 計量による実現．}
次数と一般性条件を満たす退化族を先に固定し，一つの木系 $T$ と
全ての十分小さい複素 $t\ne0$ に対する射影 $\pi_t$ を選ぶ．
各 $\pi_t$ は連続であり，$f_t=p_T\pi_t$ の全域 distortion が零へ向かう．
この選択は $t$ 方向に連続又は滑らかな fibration 族を与えるとは限らず，
任意の先指定木系・critical values への適合も主張しない．
\item \emph{作用適合 matroidal atlas．}
五次・六次では，各指定 Kähler 形式に対して atlas を持つ fibration を選べる
（定理\ref{main:action-atlas}）．しかし，これと KE 計量実現を
\emph{同じ} $\pi_t$ で同時に達成することは未証明である
（注意\ref{main:atlas-ke-scope}）．七次以上では，実際に GH 極限を実現する
$\pi_t$ の基底の各 core 近傍が位相的に非 matroidal となる
（系\ref{ke:global-core-nonmatroidal}）．
\end{enumerate}

基底は，自由な辺を持たないコンパクトな分岐二次元複体である．
「境界がない」とはこの自由辺の不存在を意味し，分岐辺の内点には三枚の sheets が接する．
その fiber は $S^1\times(S^1\vee S^1)$ であり，
core 頂点の fiber は affine complement の logarithmic Weinstein skeleton である．
全 fiber の飽和と overlap 上の partition 一致から閉同値関係を作り，
大域商の Hausdorff 性・全射性・連結 fibers を得る
（補題\ref{glue:closed-whole-fiber-quotient}）．
Cubical 比較写像 $f_t=p_T\pi_t$ の全 fiber も連結だが，
六次以上では内部木上の torus の合併に三次元部分が現れ得る
（系\ref{topo:cubical-composite-connected}）．
従って $f_t$ 自体を Lagrangian torus fibration とは呼ばない．
全平面構成の fiber Euler 関数は node と core values だけに支持され，
各 node の寄与は $1$，各 core の寄与は $\binom{d-3}{2}$ である．
有限層化に対する $\chi_c$ の加法性から，
$\chi(X_d)=d\binom d2+d\binom{d-3}{2}$ と，
$p_T$ で同一視された全 fiber 合併の Euler 数を同時に得る
（命題\ref{inv:euler-pushforward}，系\ref{inv:cubical-fiber-euler}）．
この確認は数値から全空間を同定するものではなく，
実際に構成した全域写像の各 fiber を用いる．
正則な底空間には action--angle theorem による integral affine structure がある．
core と分岐辺を含めた matroidal charts の存在，その action lattice との整合，
および既知の algebraic filling との同定は，それぞれ追加の問題として扱う．

この設定は almost toric geometry の射程とも整合する．
Leung--Symington の閉 almost toric 四次元多様体の分類
\cite[Theorem 2.1, Table 1]{LeungSymington}は，通常の曲面を基底とし，
elliptic および focus--focus singularities を持つ fibrations を対象とする．
基底の境界がない場合には K3，Enriques 曲面，torus bundles が現れる．
adjunction により $K_{X_d}=(d-4)H$ を持つ一般型超曲面については，
本稿で許す branch と core の singularities が構成の本質を担う．
action--angle coordinates，nodal trade，almost toric geometry の基本事項は
Evans の講義録\cite{Evans}に従う．

より以前にSmith\cite{IvanSmithTorus}は，
\emph{homologically essential Lagrangian torus fibers}と
\emph{tame singularities}を仮定した閉symplectic四次元多様体について，
K3または一定のtorus bundlesへの位相的制限を証明している．
この定理も今回の一般型曲面に標準的なalmost toric
singularitiesだけを使うことの困難さを示すが，
\emph{本稿の分岐fiberとcore skeletonがSmithのtame条件を満たすか，
またregular torusのhomology classがessentialかを
追加検証せずに彼の定理を適用しない}．
我々のsingular fibersは，分岐辺に沿う
$S^1\times(S^1\vee S^1)$と，
nodal points上のpinched torusに加えて，
複数のstrataを持つcore Weinstein skeletonを含む．
従って通常のalmost toric fibrationと同じ局所型を
仮定したものではない．

\subsection{Hypersurface の pants 分解と一般型曲面}

代数的 hypersurface の大域的な torus fibration を考える上で，
Mikhalkin の pair-of-pants 分解\cite{Mikhalkin}は基本的な出発点である．
同論文の Theorems 1 および $1'$ は，tropical polyhedral complex への
stratified fibration と，高次元の pair of pants による分解を与える．
基底の homotopy type は幾何種数に等しい個数の球面の wedge である．
そこでの symplectic な結論は，各最大次元 cell に少なくとも一つ
Lagrangian fiber が存在するというものである．
全ての regular fibers の Lagrangian 性を得るためには，
接着に用いる form と local fibrations の整合をさらに示す必要がある．

一般型 hypersurface の Calabi--Yau 的な部分領域は，
Leung--Wan\cite{LeungWan}によって研究されている．
同論文は holomorphic volume forms を局在させる分解を構成し，
Mikhalkin の構成で既に Lagrangian である fibers について
asymptotically special となる性質を示す．
一般型における SYZ 的構造を調べる動機として，
この解析的な結果と，全 regular fibers を Lagrangian にする存在問題の両方が重要である．

Zhang\cite[Theorem 1.1]{Zhang}は，次数 $d$ の超曲面と微分同相な四次元多様体を，
$d(d-4)^2$ 個の四次元 pants と K3 に由来する領域などに分解している．
その pants 領域の境界には，torus を一般 fiber とする graph への fibration が構成される．
参照した preprint の Figure~1.1 と Theorem~1.1(ii) における特異 fiber は
$\Theta\times S^1$ であり，$\Theta$ は二つの三価頂点を持つ theta graph である．
これは本稿の一つの四価頂点を持つ $(S^1\vee S^1)\times S^1$ と同相ではない．
同論文の Introduction には，
五次曲面の五平面退化から五個の pants を得る記述もある．
本稿の五次の場合は，この既存の幾何学的構図を踏まえている．
一方，torus に対する積分 $\int_T\omega=0$ と，
各 tangent space 上での条件 $\omega|_T=0$ は異なるため，
本稿では後者を各 piece とその overlaps で直接確認する．

Golla--Martelli\cite{GollaMartelli}は，四次元多様体の pants 分解を
微分位相的な構成法として展開し，corner を含む貼合せと mapping-class data を記述している．
本稿でも，分岐底空間の cells と boundary faces の構造を理解する際にこの記述を用いる．
symplectic な構成では，これに fixed collars と実際の action periods の整合が加わる．
境界を固定した core の置換を検討する際には，corner に追加の torus 領域を置く操作が
基本群と boundary topology を変え得ることも考慮する必要がある．

\subsection{SNC divisor の有限相対 compactification}

Evans--Mauri\cite{EM}は，contact link 上の Lagrangian fibration を
Liouville flow によって completion 上へ延長する方法を与えた．
一般位置の四直線を $\PP^2$ から除いた四次元 pants に対しては，
$K_4$ の cone を基底とし，torus，$S^1\times(S^1\vee S^1)$，
core skeleton を fibers とする模型を構成している
\cite[Lemmas 2.16, 2.22; \S3.1]{EM}．
同論文は一般型曲面への応用も明示的に動機として挙げている．
また Remark 2.23 は，曲線上の Lagrangian leaves が一次元であるため，
circle-bundle lift に必要な curvature の vanishing condition が自動的であることを指摘する．
この点に genus 零の制約はなく，本稿の positive-genus 構成も同じ lifting principle を用いる．

本稿の相対問題では，completion の内部への cone extension に加え，
所定の divisor における有限時刻の circle collapse を制御する．
Lemma 2.16 の completion は contact form の正の関数倍を許すため，
固定した finite collar，normal moment，divisor 上の foliation，
core 近傍の primitive を同時に指定するには追加の構成が必要である．
同じ contact 構造を保つ正の rescaling でも既存の葉の Lagrangian 性は一般に変わる
（注意\ref{ke:contact-rescaling-collar-scope}）．
一方，\cite[\S\S2.3--2.4]{EM}の Hermitian metric の変更は選んだ Kähler class を，
exact symplectic deformation は de Rham class を保つので，
固定 cohomology class 自体を本稿固有の特徴とはしない．
ここでの相対条件は，Kähler representative と指定 product germs，
marked full leaves，選択する holonomy と periods，有限 endpoint の同時制御である．
SNC symplectic neighborhoods とその relative regularizations には
Tehrani--McLean--Zinger\cite[Definitions 2.11--2.12, Theorem 2.13]{TMZ}の枠組みを用いる．
同定理は一般には symplectic form の変形を許すため，本稿の固定 $\omega'$ に対する
charts は，product crossing germs を準備した後の相対 tubular Moser で構成する．
正の weighted divisor に適合する primitive の存在は
Borman--Sheridan--Varolgunes\cite[Propositions 2.7, 2.10]{BSV}に関係する．
同結果の重みは正の実数でよく，$H^1(M)=0$ を要求しないが，先指定した normal connection，
core primitive germ，full Morse leaves の同時保存までは結論していない．
同論文の period-one 規約とは $r_a=2\pi\mu_a$，$\kappa_a=2\pi w_a$ で対応する．
normal directions の対称化は Pascaleff\cite[\S4.2]{Pascaleff}の構成と関係する．
Kähler class の実数係数の表示から global logarithmic potential を作る段階では，
compact Kähler manifold の $\dd\dd^c$-lemma\cite[\S5]{DGMS}を用いる．

具体的には，$M$ を射影曲面，$D=\bigcup_aD_a$ を三重交点のない SNC divisor とし，
各 $D_a$ が少なくとも二つの intersection points を持つとする．
$D_a$ の genus，self-intersection，同じ二成分の複数交点には制限を加えず，
\[
 [\omega]=2\pi\sum_a w_a[D_a],\qquad w_a>0
\]
を仮定する．
定理\ref{snc:main}は，同じ class の Kähler representative $\omega'$ に対し，
$D$ を outer graph の逆像とする closed-cone fibration を与える．
各 punctured component 上には，あらかじめ指定した無 extrema の
marked Morse foliation を用いる．
full saddle leaves と finite endpoint を保ち，
交点の product germs および有限個の divisor-compatible flat patches も指定できる．
個々の divisor components の ampleness や $H^1(M;\RR)=0$ は仮定しない．

positive genus における period の制御には，connection holonomy の自由度を用いる．
固定 residues を持つ二つの collar primitives の差は，
plumbing の homology exact sequence により $H^1(D;\RR)$ に由来する．
これは dual graph の periods のほかに，各 $D_a$ の $2g_a$ 個の periods を含む．
crossing disks で零の closed one-form $\eta_a$ により normal connection を変え，
対応する coupling forms を Moser map $\Phi_s$ で固定すると，
\[
 [\Phi_s^*\lambda_{a,s}-\lambda_{a,0}]
       =-s w_a[\eta_a]
\]
となる（補題\ref{snc:holonomy}）．
これで component periods を消し，残る dual-graph periods を basic forms で補正する．
この Moser 操作は auxiliary normal-bundle charts を変えるものであり，
actual collar 上の Kähler form $\omega'$ は固定される．
最後に bounded primitive の cutoff を用いて，
core 近傍の primitive germ を保ったまま end の model と接続する．
さらに connection と coupling charts を固定し，closed basic correction と exact cutoff
に限る場合の必要十分条件を命題\ref{snc:fixed-holonomy}で与える．
Reeb 商 $q:D\to\Gamma$ に対し，障害は
\[
 H^1(D;\RR)/q^*H^1(\Gamma;\RR)
 \cong\bigoplus_a\operatorname{Hom}(L_a,\RR),\qquad
 \operatorname{rank}L_a=g_a
\]
であり，各成分から $g_a$ 本ずつ選んだ正則葉の period の消滅で判定できる．
曲面一枚での leaf-period と basic representative の関係は
Izosimov--Khesin--Mousavi\cite[Lemma 4.8, \S5]{IKM}の既知結果である．
本稿では整数格子の証明と SNC 全体の exact sequence を記述し，有限時刻の全葉 collapse と
core に相対的な接着へ適用する．任意の障害の消去には $\sum_a g_a$ 個の独立な holonomy
調整方向が必要十分である（系\ref{snc:minimal-holonomy-controls}）．
非零障害を消す connection 変更は一般に全 full saddle collars を支持外にできない
（注意\ref{snc:saddle-support-obstruction}）．これらは指定した補正方法の条件であり，
任意の Lagrangian fibration の存在を分類するものではない．

固定した元の form が交点で nonorthogonal な場合には，
commuting normal-circle actions の product model をそのまま要求できない．
これは\cite[Proposition 2.5]{BSV}の orthogonality 条件であり，
本稿では Kähler class 内の明示的な potential deformation によって product germs を作る．
さらに global residues と core の primitive germ の関係を記述し，
collar の area balance だけから global exactness を推論しない．
$\PP^1\times\PP^1$ の toric boundary，
二点で交わる二本の $(1,1)$ curves，
九点で交わる二本の smooth plane cubics は，
nonample components，multiple intersections，positive genus をそれぞれ検査する例になる．

\subsection{実際の退化族上での接着}

Fernández de Bobadilla--Pełka\cite[Proposition 5.6, Theorem 6.14]{FBP}は，
Kähler SNC 退化の近傍で ivy-like な Lagrangian torus fibrations を構成し，
一般型を含む広い範囲の退化を扱う．
曲線 stratum 上の Morse saddle から branch が生じる点は本稿にも共通する．
同論文が大域的に構成するのは，maximal および submaximal strata に対応する，
nearby fiber の境界・角を持つ余次元零の部分領域である．
component の内部に対応する残余領域を埋めることは追加の問題であり，
同論文の Example 5.8 もこの差を説明している．
本稿では，この残余領域に core を含む finite compactification を入れ，
double-curve necks および triple-point models と有限の overlaps 上で接着する．
同論文の Theorem~6.14(b) の fibration の定義域はこの部分領域だが，
同定理(c) の GH 収束は構成した Kähler 計量を入れた $W_z$ 全体を扱い，
退化写像が proper なら $W=X$ と取れる．従って同論文の GH 結論まで
部分領域だけのものとは解釈しない．その計量と本稿の actual 負 KE 計量の区別は
第\ref{intro:ke-collapse}節で述べる．

central fiber から symplectic parallel transport で fibration を移す考え方には，
Ruan の Calabi--Yau hypersurface の構成\cite{Ruan}がある．
Ruan の五次は $\PP^4$ 内の Calabi--Yau 三次元多様体であり，
本稿の $\PP^3$ 内の一般型五次曲面とは異なる次元の問題を扱う．
また Mak--Matessi--Ruddat--Zharkov の共同研究について，
Ruddat の二〇二六年六月の講演概要\cite{RuddatAnnouncement2026}では，
reflexive polytope に対応する toric Fano manifold 内の compact Calabi--Yau hypersurface 上で，
余次元二の discriminant を持つ Lagrangian torus fibration の構成が発表されている．
概要は Evans--Mauri の方法に従う Liouville flow と帰納的な mixed local models を挙げる．
これは関連する大域的構成の動向として引用し，本稿での比較は公表概要の射程にとどめる．
本稿の一般型曲面の SNC cores と整数分母の問題には，指定した局所模型の接着条件が加わる．
さらに Zharkov の二〇二六年四月の講演概要\cite{ZharkovSoftCells2026}は，
phase-tropical および ober-tropical な超曲面の比較から，
permutahedral skeleton を soft cells に置き換える局所模型を述べている．
元の超曲面との symplectic な比較と Lagrangian skeleton fibers が発表の対象であり，
pants の局所構成を比較する上で関連する．
これも公開された講演概要として引用し，指定した SNC data に相対的な接着定理を
その概要だけから導いたり，完成論文の証明を確認したものとして扱ったりしない．

四次曲面における transport と moment 保存の記述は
Evans\cite[Chapter 6]{EvansNotes}にも詳しい．
Chakravarthy--Groman\cite[Theorem 1.6]{CGK3}は，
symplectic Kulikov model の component 上の almost toric fibrations を
nearby K3 へ接着する大域定理を証明している．
その component hypotheses には anticanonical cycles が含まれるため，
多数の punctures と saddles を持つ本稿の曲線 strata には別の collar 構成が必要になる．
同論文の改訂版 Lemma 2.13 が扱うのは，必要な cylinders を潰した
central fiber の商への transport の連続性である．
本稿では smooth strata 内の compact saturated subsets を運び，
double curves と triple points の中心部分を正確な $xy=t$，$uvw=t$ の models で覆う．

位相型と linear monodromy だけでは symplectic gluing は決まらない．
曲面上の Morse systems には weighted または measured Reeb graph の情報があり
\cite{DMT,IKM}，compact focus--focus fiber の近傍には
semi-global invariant がある\cite[Theorem 2.1]{VuNgoc}．
一重 focus の十分小さい飽和近傍では $H_1\cong\ZZ$ であり，
正則 torus の normal vanishing cycle は消え，残余 cycle が generator となる．
従って既に存在する全近傍の symplectic map が指定 primitive に関して exact であることは，
その一つの残余 period の一致と同値である（補題\ref{ke:focus-neighborhood-period}）．
Node 上の primitive germ だけを固定してもこの period は変えられるが，
一つの外側正則 fibre 上の primitive まで固定すれば障害は零となる
（例\ref{ke:focus-local-primitive-shift}）．これは constant-zero Taylor 不変量と，
\cite[\S7]{VuNgoc}の exact 版に現れる作用の定数項との区別を具体化する．
そこで，本稿は actual family 上の同じ Hamiltonian circle action と
normal moment $\mu$ を用い，reduced surface 上の functions $h_\mu$ を
marked foliation に相対的に補間する．
$\{\mu,h_\mu\}=0$ によって Lagrangian 性を保ち，
既に構成した nodal fiber はその compact fiber 全体と近傍を transport する．
その際に葉の分割と Eliasson 座標・cycles を同時に運べば Taylor 不変量が保存され，
exact 移送では作用の定数項も保存される（命題\ref{ke:focus-whole-neighborhood-transport}）．
この方法は測度や period の不変量を無視するものではなく，
同じ実際の form 上で leaves と action differentials が一致するように作るものである．
正則 torus band の既定 base 同定については，normal moment の一致に加え，
第二作用が整数 shear と定数 translation を除いて一致することが必要十分である
（命題\ref{neck:regular-period-matching}）．固定 primitive に関する exact 同定では
その translation も零を要求する．これは action--angle theorem の局所的な帰結であり，
独立な第二作用の自動一致や，任意の先指定 collar map の大域延長を主張しない．
circle-invariant nodal fibrations の reduction による記述には
Groman--Varolgunes\cite[\S7.3]{GV}を用いる．

component 上の局所点 blow-up には，
Chakravarthy--Groman\cite[Theorem 1.1]{CGblowup}の almost toric comparison と，
Fleming\cite[Proposition 4.6, Construction 4.7]{Fleming}の Kähler blow-up を用いる．
これらは regular elliptic-boundary patches 内で適用し，
core と crossing germs を固定する．
容量は例外曲線の実面積 $e$ で測り，実際の複素 blow-up 上の radial K\"ahler form と
半径 $\sqrt{e/\pi}$ の symplectic blow-up の相対同定を直接構成する
（補題\ref{glue:area-normalized-blowup}）．先に飽和した affine-convex patch を取り，
その内部へ球と容量 triangle が入るよう $e$ を小さくする．比較写像の恒等域はその
飽和 patch の外であり，球の外全体ではない．
その後の divisor-relative Moser と normal-circle conjugacy は，
曲線の genus に依存しない．
positive genus の場合も，同じ実際の marked curve 上の isotopies を用い，
normal-bundle の degree の一致に加えて line-bundle の smoothing relation を保持する．

\subsection{次数分割，偏極，底空間の整数位相}

まず平面退化の場合，定理\ref{main:global}は $d$ 個の cores，
$d\binom d2$ 個の nodal fibers，
$\binom d2(d-4)$ 本の branch edges を持つ fibration を構成する．
$d=5$ の底空間は四単体の二次骨格である．
一般の $d$ では，同じ実際の double curves 上の marked pants trees から
有限複体 $B_d(T)$ を作り，
\[
 B_d(T)\simeq\bigvee^{\binom{d-1}{3}}S^2
\]
への整数鎖ホモトピーと $H_2$ の具体的な基底を与える．
球面の wedge という型は Mikhalkin の結果と整合するが，
ここでは指定された branch edges，bigons，triangles に対する鎖データを記述する．

任意の順序付き分割 $d=e_1+\cdots+e_k$，$k\geq3$ に対しても，
一般の smooth hypersurfaces $S_i=\{F_i=0\}$ を用い，
$\prod_iF_i+tG_d=0$ を順序付き小解消することができる．
定理\ref{part:global}は，この退化から $k$ 個の cores を持つ fibration を与える．
第 $i$ core は
$S_i\setminus\bigcup_{j\ne i}(S_i\cap S_j)$ の logarithmic skeleton に由来し，
node 総数は $d\sum_{i<j}e_i e_j$，
branch edge 数は $(d-4)\sum_{i<j}e_i e_j$ である．
定理\ref{part:homotopy}はその底空間について
\[
 B_{\mathbf e}(T)\simeq
 \bigvee^{\,p_g(X_d)-\sum_i p_g(S_i)}S^2
\]
を示し，Reeb graphs の cycles，同じ triple-point type の複数交点，
および tetrahedron boundaries から整数基底を構成する．
例えば次数六の分割 $(2,2,2)$ と $(1,1,4)$ は，
同じ symplectic 曲面上に $b_2(B)=10$ と $9$ の底空間を与える．
従って，core の選択は底空間の homotopy type にも反映される．

偏極の範囲はこの自由度とは別に正確に記述できる．
小解消の vertical class を
$\xi=aH+\sum_i\beta_i[D_i]$ とし，
$\delta_\ell=\beta_{\ell+1}-\beta_\ell$，
$E_\ell=e_1+\cdots+e_\ell$ と置く．
定理\ref{part:cone}は，その Kähler slice が
\[
 \delta_\ell>0\quad(1\leq\ell<k),\qquad
 \sum_{\ell=1}^{k-1}E_\ell\delta_\ell<a
\]
という sharp な開単体であることを示す．
平面分割では $E_\ell=\ell$ となり，定理\ref{pol:cone}を回収する．
この slice 全域の記述には Bézout と Nakai--Moishezon を用いる．
大域 fibration の局所 blow-up construction に使うのは，
一様な patches を確保した後の十分小さい ray である．

\subsection{core の topology と matroidal charts}

core の homotopy type と embedded Liouville skeleton の同定は別の問題である．
Hattori\cite{Hattori}は一般位置の hyperplane complement の homotopy type を決定し，
関連する arrangement の記述は Papadima--Suciu\cite[\S1.1]{PapadimaSuciu}にもある．
Ruddat--Sibilla--Treumann--Zaslow\cite{RSTZ}の combinatorial skeleton は，
Zhou\cite{Zhou}によって，適切な exact symplectic structure の
Liouville skeleton として実現されている．
本稿では，指定された compactification の collar を保つことが必要なため，
実際の logarithmic potential から得られる skeleton を用いる．
特定の RSTZ model との相対同定は追加条件として残る．
Matessi\cite{MatessiPants,Matessi}の Lagrangian pants は半次元の部分多様体を作る構成であり，
symplectic pants 全体の torus fibration とは区別される．
Abouzaid--Auroux--Katzarkov\cite{AAK}の対象も，
hypersurface に付随する toric blow-up や conic bundle であり，
compact hypersurface 自身への本稿の gluing theorem は別途必要である．

一般型曲面内の部分的な almost toric geometry にも先行研究がある．
Evans--Urz\'ua\cite{EvansUrzua}は，symplectic form を変形した一般型曲面で，
rational homology balls の非有界な列の symplectic embeddings を構成した．
Brendel--Hauber--Schmitz\cite{BHS}は，$B_{dpq}$ の skeleton 近傍からの
partial embeddings によって移される小さい Lagrangian tori を
ambient symplectomorphism に関して区別し，一般型の五次曲面や
simply connected Godeaux surface にも，適切な symplectic form の下で応用している．
ここでは displacement energy germ が重要な不変量である．
本稿が指定する K\"ahler class，compact な底の全体，core と neck の同時接着は
別途確認する条件であり，局所 torus の非同値性から fibration 全体の非同値性を推論しない．
この区別は，構成後の torus fibers や atlas の同値性を研究する際にも維持する．

torus fibers の Floer 不変量を比較する方向では，Venugopalan--Woodward
\cite[Definitions~1.1--1.2, Theorem~1.3]{VenugopalanWoodward2026}が，
compact almost toric 四次元多様体の torus fiber に対する Maslov index 二の disk counts を
tropical graphs で表す公式を与えている．同定理は elementary polyhedral decomposition と，
potential に寄与する全 disks の tropical graphs の collisions が generic interior points にある
という仮定を伴う．本稿の分岐底複体と core fibers はこの almost toric 設定に含まれず，
公式を本稿の fibers に用いるには別の比較定理が必要である．
また非単調な場合の potential は Maurer--Cartan space 上の関数を同型まで考えるものであり，
選択に依存しない一つの Laurent polynomial とみなして atlas の非同値性を判定しない．

matroidal charts との比較には，Bergman complex の位相を使う．
Ardila--Klivans\cite{ArdilaKlivans}は lattice of flats による subdivision と
phylogenetic-tree spaces との関係を示し，
Feichtner--Sturmfels\cite[Theorem 4.1]{FeichtnerSturmfels}は
nested-set complexes による unimodular triangulation を与えた．
Petersen graph と $K_{3,3}$ は後者の Examples 1.1--1.2 に既に現れる．
Shaw\cite[\S2.4]{Shaw}の matroidal fans は，
本稿で比較する smooth tropical local models の範疇を与える．
これらの既知の模型を含め，定理\ref{trop:cubic-classification}では，
二価点を除いたリンクが三価である rank-three matroidal fans を分類する．
連結 matroid から現れるのは
$K_4,K_{3,3},\mathrm{Petersen},\mathrm{Heawood}$ であり，
非連結 matroids と Euclidean factors を含めてもリンクの第一 Betti 数は八以下である．
体上での realizability はこの上界の仮定ではない．

曲線配置の tropicalization を matroid の外側で扱う比較対象として，
Aksnes\cite[Theorem 4.15]{Aksnes}は，平面内の transverse very affine arrangement の
補集合を arroid fan によって記述する．
本稿の Reeb 展開を施した contact link と，divisorial tropicalization の support が
同相であるとは限らない．また，この結果から actual action periods の一致は従わない．
従って，matroidal charts の障害の後に比較する，より広い fan の範疇として区別する．

一方，平面退化の各 core link は
$b_1=\binom{d-2}{2}$ を持つ．
従って $d\geq7$ では，その core の局所位相だけで matroidal chart を排除できる
（定理\ref{trop:degree-classification}）．
次数分割の場合も
\[
 b_1(\Gamma_i)=1+\frac{e_i(d-e_i)(d-4)}2
\]
であり，同じ障害が全ての partitions に及ぶ
（命題\ref{trop:partition-obstruction}）．
六次では，全ての core が matroidal な局所位相を持つ候補は全平面退化に限られる．
その全 core links が Petersen となる labelled split systems は正確に六個で，
$K_6$ の one-factorizations と対応する（定理\ref{topo:petersen}）．

五次では Betti 数だけの判定をさらに精密化する。
divisor の各 punctured component の内部 Reeb graph は、core link の誘導部分グラフになる。
この block と個々の triple-point labels を保持すると、$(1,1,3)$ と $(1,2,2)$ の
全 core の同時 matroidal 化は不可能である一方、$(1,1,1,2)$ では四つの標識付き
abstract graph systems が存在する（定理\ref{q5:classification}）。
その条件は三つの quartet splits の parity で記述できる。
Reeb graph の実現には good orientation が必要であり、
Michalak\cite[Definition 2.2, Theorem 5.6]{MichalakReeb}の実現問題との関係を明示した上で、
本稿で使う境界付き曲面には saddle pants と cylinders による直接の構成を与える。
従って有限グラフの個数を、embedded marked isotopy classes の個数とは同一視しない。


\subsection{局所作用座標から大域 atlas へ}

Petersen link が得られても，標準 fan の格子と実際の作用格子を合わせる必要がある．
critical slopes の rationality，annular area の Farey 条件，branch lattice の saturation は，
そのための局所的な条件である（命題~\ref{trop:farey-condition}，\ref{trop:branch-denominator}）．
固定 annular area は $\operatorname{area}/(2\pi w)=1/m$，$m$ は正の奇整数，という
離散条件を持つ（命題~\ref{per:odd-reciprocal}）．
一方，marked isotopy class 内で foliation の代表を選べるなら，五直線補集合で各 critical
basin の面積を $2\pi w_b$ にし，内部 annulus の面積を $2\pi w_a$ にできる．
これにより十五面全ての作用格子が標準 Petersen fan の格子と一致する
（定理~\ref{per:local-realization}）．
これは Evans--Mauri の Liouville cone construction\cite{EM}と，標準 tropical lattice
\cite{GKM}を，実際の作用積分によって接続する構成である．
四直線の $K_4$ core については，有向四周期に沿う basin areas を用い，六面の作用座標と
$H_1(W_4;\ZZ)\cong N_4^*$ の整数同型を与える（命題~\ref{per:k4-realization}）．
Petersen の全 meridians に一つの ambient-linear extension を要求すると半整数の障害があるが，
この追加条件は一つの core の作用適合 chart の存在とは異なる．

大域 atlas では，さらに chart transitions を検査する必要がある．
Shaw\cite[Definition~3.1]{ShawSurfaces}，Jell--Rau--Shaw
\cite[Definitions~2.1, 3.1, 5.1]{JellRauShaw}，L\'opez de Medrano--Rinc\'on--Shaw
\cite[Definitions~2.2--2.4]{LRSChern}の枠組みでは，transitions は ambient integer affine maps
の制限である．face ごとの unimodularity と branch saturation だけでは足りない．
本稿は三価 star の唯一の線形関係を用いて延長条件を記述し，両 chart で同じ Hamiltonian
circle の作用が ambient-affine ならその条件が成立することを示す
（補題~\ref{atl:branch-extension}，命題~\ref{atl:common-hamiltonian}）．
focus--focus points では作用格子の monodromy が非自明である\cite[Remark~3.1]{VuNgoc}ため，
atlas はその有限集合を除いた底空間に与え，除いた点には実際の focus affine singularities を残す．

Evans--Mauri の掲載版 Propositions~6.11--6.12 と Theorem~6.14 は，
Assumption~1 の下で toric neighborhoods と双対 monodromy を比較する\cite{EM}．
この仮定は各一次元 stratum が $\PP^1$ で零次元 stratum と正確に二点で交わることを要求し，
各直線に $m-1\ge3$ 個の交点を持つ本稿の一般位置 $m$ 直線配置には直接適用できない．
本稿の三穴・四穴 sphere とその saddle branches には，critical actions を直線に保つ別の接着が必要である．
実際，作用適合 matroidal chart が存在すれば，各 incident sheet の critical action は
normal moment の有理勾配の affine function でなければならない
（命題\ref{atl:critical-action-obstruction}）．
固定した exact symplectic neck 上で，normal action，両端の open germs，toric collars，
saddle 点の局所 germ，marked tree を全て保ちながら，critical leaf の regular arc だけを動かして
この作用を非直線にできる（例\ref{atl:fixed-form-bent-branch}）．
その点には任意の分母を許しても作用適合 matroidal chart はなく，
Lagrangian 性と atlas の存在が別の条件であることを示す．
このため，actual neck の全空間上で一つの primitive を取り，三重点の toric reference actions
と両側の作用の加法定数を同時に固定する．
固定した正負の normal moments における symplectic transport はこの primitive を exact 差まで保つ．
境界作用に対する Stokes の公式から reduced total area の affine 性を得て，
critical basins と内部 annuli の面積を指定する parameter に相対的な Moser construction を行う．
これにより有限 $t\ne0$ でも critical action functions を直線に保てる
（系~\ref{ng8:central-criterion}）．
measured Reeb graphs\cite{DMT,IKM}との関係では，既に固定した二つの systems を
無条件に同一視するのではなく，同じ actual reduced forms 上で必要な critical-area data を実現している．

この接着を用いると，五次と六次の全平面退化で大域 atlas を得る
（定理~\ref{ga8:global}）．$b_i=\beta_i/(2\pi)$ とし，
\[
 c=\frac{a}{2\pi(d-1)},\qquad w_{i,j}=c+b_j-b_i
\]
という weights を選ぶ．五次では $f_i(j)=(i+j)/2\bmod5$ の対称な四周期系を用い，
各共有曲線の五つの blow-up points を三つの basins へ $(3,1,1)$ と配分する．
六次では各曲線の四 basins へ一つずつ，内部 annulus へ二つを配分する．
面積変化と normal framing の変化が同じ整数条件を与え，両端の affine actions が一致する．
五次の各 triangle には五つの nodes，六次の各 triangle と bigon にはそれぞれ三つと二つの nodes が入る．
全 branch lengths は $a/(\pi(d-1))$ となる．
五次の対称な標識付き recipes と六次の Petersen systems はそれぞれ六個あるが，
この数を非同型な大域 atlases の個数とは同一視しない．

\subsection{混合分割での局所存在と大域障害}

分割 $(1,1,1,2)$ では，quadric の三本の $(1,1)$ curves の補集合に
$K_{3,3}$ chart を構成する（定理~\ref{quad:realization}）．
標準 $K_{3,3}$ fan が二つの tropical lines の積であることは，parallel connection に対する
Shaw--Werner\cite[Theorem~7.2]{ShawWerner}の記述に従う．
ここで必要な追加の計算は，自己交点二の actual normal framing と非零 critical periods である．
各 component の一つの basin をその weight 分だけ広げると，内部 annulus の面積比が一となり，
九面の作用行列は全て unimodular となる．
二直線と conic を境界に持つ三つの plane cores にも，同様に非零 periods を用いた
$K_4$ charts がある（命題~\ref{mixper:realization}）．
従って，四つの標識付き quartet systems の全てで，各 core を個別に実現できる．

分岐における作用の不整合は，tropical divisor の言葉で表せる．
単位 facet weights を持つ三価 star の primitive transverse lifts を
$u_1,u_2,u_3$，分岐方向を $v$ とし，
$u_1+u_2+u_3=\kappa v$ と書くと，sheet ごとに整数アフィンな作用関数 $h$ の係数は
\[
 \mathfrak d(h)=\sum_{a=1}^3\dd h_a(u_a)-\kappa\dd h(v)
\]
となる．これは Allermann--Rau の tropical principal divisor の係数そのものである
\cite[Definition~3.4, Remark~3.5]{AllermannRau}．
同文献の abstract cycles に対する構成
\cite[Definition~6.1, Construction~6.4]{AllermannRau}と整合するが，
本稿では各 divisor component に固有の normal action を用いるため，
全 core の係数を初めから一つの大域 principal divisor と解釈しない．
共有 neck の normal circles が互いに逆向きであることから，
対応する二端の係数が相殺するという必要条件を直接用いる．

分母の自由度を除くのは，局所 fan の分類だけではなく，
共有する四 puncture 曲線の annulus と nodal framing の整合性である．
一つの focus node が与える integral shear は almost toric geometry の標準的な事実である
\cite[\S4.3, equation~(4.10)]{SymingtonFour}．
本稿では，両側の annulus の unimodularity と四 punctures の整数 cycle shears を同時に用い，
$q_Aq_B\mid2$ を導く（命題~\ref{q9:rigidity}）．SNC の node budget を合わせると，
二端の分母 $q_A,q_B$ に対して内部 annulus の node 数が
\[
 n_{\rm int}=\frac{2}{q_Aq_B}
\]
に等しいことを示す．三価 branch の saturation により分母は奇数であるから，
node 数の整数性は $q_A=q_B=1$ および $n_{\rm int}=2$ を強制する．
この量子化条件を混合五次の三本の conic necks に適用すると，
高分母を許した場合にも core coefficients の総和は $-6$ となる．
従って現在の SNC cores と branch graph を保つ範囲では，
混合五次の大域作用適合 matroidal atlas は分母の選択によらず存在しない．
これは regular curves と annuli の周囲の議論であり，
focus values 自体に matroidal chart を延長する主張ではない．

三本の plane--plane branches の分母を任意のまま残しても，line components の寄与は相殺される．
各 mixed plane core の係数総和は $-1$，quadric core は $-3$ なので，全体は $-6$ となり，
対向端の相殺条件に反する（定理~\ref{hd9:mixed-no-go}）．
これと位相障害を合わせると，現在の比例曲率 SNC 構成で cores と branch graph を保つ
大域作用適合 matroidal atlas が可能な次数分割は，分母を制限せずに，全平面の五次と六次に限られる
（系~\ref{hd9:classification}）．

同じ四 puncture の剛性は，六次の全三十 branches の分母が一であり，
十五 bigons の各々に二 nodes が入ることも強制する．これは一つの weight recipe に依存しない．
一方，四直線の局所 $K_4$ core では，正の basin areas と六つの unimodular action matrices を
同時に指定することで，任意の奇数 $q\ge3$ に対して分母 $(q,1,1,1)$ を持つ chart を得る
（\cite{QuinticArithmetic}）．この例は，局所の奇数分母条件と，実際の共有 neck に沿う
大域的な剛性との違いを具体的に示す．

\subsection{負 Kähler--Einstein 退化との計量的な比較}\label{intro:ke-collapse}

定理\ref{ke:global-realization}は，Lagrangian fibration の射影に標識付き基底収縮を合成した写像によって，
負 KE 計量の径の尺度での Gromov--Hausdorff 極限を実現する．
Actual KE 形式に関する fibration $\pi_t:X_t\to B_d(T)$ と
標識付き収縮 $p_T$ の合成について，全曲面の ambient 距離に対する
一様な距離歪みが消えることを証明する．
正規化は $\operatorname{Ric}(\omega_t)=-\omega_t$，
$\bar g_t=R_t^{-2}g_t$，$R_t=\log(-\log|t|^2)$ である．
以下では既存結果の範囲，本稿が明記して用いる局所入力，
および全域化のために証明する評価を区別する．

本稿の fibration の存在と，負 Kähler--Einstein 計量の崩壊は異なる主張である．
全平面退化 $X_t=\{\prod_i L_i+tF_d=0\}$ に対し，
$\operatorname{Ric}(\omega_t)=-\omega_t$ と規格化すると
$[\omega_t]=2\pi(d-4)H$ なので，本稿の存在定理の偏極範囲に入る．
しかし，各 $t$ で $\omega_t$-Lagrangian な fibration が存在することは，
その fibers の KE 直径が小さいことや，退化 parameter に関する計量的一様性を含意しない．

Song\cite[Theorem~1.2]{SongNegativeKE}の stable degeneration の定理を
この族に適用すると，成分内部に基点を選んだ非崩壊極限は
$\PP^2$ から $d-1$ 本の直線を除いた $d$ 個の完備 KE 曲面となる．
従って元の尺度では，各成分の内部領域には四実次元幾何が残る．
固定成分の空間的 blow-down は companion 原稿の Part I
\cite[\texttt{cor:arrangement-cone}]{MultiscaleKECurrent}の対象である．
そこでの $C_{2,d-1}$ は，本稿の compact な極限 $\mathcal C_d$ の成分頂点における
metric tangent cone と labels を含めて一致する
（命題\ref{topo:component-cone-germ}）．さらに本稿の同じ $f_t$ の distortion を使うと，
元の尺度と径尺度の間に適切な $1\ll a_t\ll R_t$ を選び，
成分頂点上の基点に関する $(X_t,a_t^{-1}d_{g_t})$ の pointed GH 極限も
この錐となる（系\ref{ke:intermediate-component-cone}）．さらに同系では，
基点の target 値の成分頂点からのずれも尺度の選択に含めることで，
component core と変更 collar の全領域内の任意の基点列へ結論を拡張する．
さらに命題\ref{topo:cubical-local-models}と\ref{topo:cubical-moving-blowups}は，
cubical target の全局所模型と移動基点の拡大極限を分類する．
$\epsilon_t=\operatorname{dis}(f_t)$ とし，$1\ll a_t\ll R_t$，
$R_t\epsilon_t=o(a_t)$ を満たす尺度では，actual 曲面の全ての移動基点の
部分列極限も，有限基点を持つ $C_{2,d-1}$，$\mathbb R\times\mathcal Y_{d-2}$，
全角 $3\pi/2$ の錐，および $\mathbb R^2$ に尽くされ，各型を実現できる
（系\ref{ke:intermediate-moving-classification}）．
ここで $\mathcal Y_{d-2}$ は原点で貼った $d-2$ 本の半直線である．
これは全域歪みの閉球への移送と有限複体の幾何による帰結であり，
上の誤差条件を除いた任意の中間尺度や極限の交換，tensor 収束は主張しない．
固定補空間の Part I--III，退化族の局所一様解析と条件付き大域枠組の Part IV--V，
本稿の actual fibration による直接実現は，この対象と量化を区別して接続する．
Part II の標識付き残余複素曲線の period と，本稿の第二作用 period も別のデータである．

他方，Zhang\cite[Theorems~1.1, 2.4]{ZhangNegativeKE}は，
全空間が $\mathbb Q$-factorial な simple toroidal canonically polarized degenerationで，
深い stratumへ基点を動かした局所 torus collapse を証明した．
曲面の二重曲線では一つの円，三重点では二次元トーラスが潰れる．
これは compact な全基底への GH 収束とは異なる．
本稿の族には $xy=tz$ 型の非 $\mathbb Q$-factorial な全空間特異点があり，
小解消後も crepant exceptional curve 上で相対 canonical bundle の次数が零となるため，
同定理の相対 ample 仮定をそのまま使えない．
Mandel\cite{MandelNegativeKE}の負 KE neck の構成も，
我々の compact family との全域的な計量比較を自動的には与えない．

Fu--Hein--Jiang\cite[\S1.1, Main Theorem]{FuHeinJiangCusp}は，
孤立楕円特異点へ退化する六次曲面の actual 負 KE 計量を接着により解析し，
形成される complex hyperbolic cusp の先端に Tian--Yau gravitational instanton が
bubble として現れることを証明している．これは負 KE 退化の実際の計量を扱う
先行結果であり，全平面 SNC 退化の全曲面を同じ Lagrangian 射影で比較する
本稿の主張とは，退化模型と極限の尺度を区別して比較する．

Fernández de Bobadilla--Pełka\cite[Theorem 6.14(b),(c)]{FBP}では，
部分領域上の ivy-like fibration と，proper な場合の全 nearby fiber の GH 収束が得られる．
その計量は，選んだ Kähler 形式への exact な補正から作る $\omega_q^\varepsilon$ による．
Proposition~6.13 と Theorem~6.14(c) はこの計量を再尺度化して双対複体へ収束させる．
本稿との比較では，領域の全域性だけでなく，二形式・尺度・射影を同時に区別する．
Calabi--Yau の場合にも，これと Ricci-flat 形式は区別されている
\cite[Introduction, Remark 7.4]{FBP}．同じ cohomology class の形式を各固定 $t$ で
Moser 移送できても，その移送の計量的な一様性は別に必要である．
実際，$\phi_t^*\omega_t'=\omega_t$ から従うのは
$\phi_t^*g_t'=\omega_t(\,\cdot\,,\phi_t^*J\,\cdot\,)$ であり，
これを $g_t=\omega_t(\,\cdot\,,J\,\cdot\,)$ と比較する評価は自動ではない．
本稿では局所移送の変位，全域 scalar 下界と portal 上界を直接評価し，
指定された actual 負 KE 計量と同じ $f_t=p_T\pi_t$ を比較する．

局所 torus fibration により近い先行結果は，
Ruan\cite[Proposition~3.1, Lemma~3.3, Theorem~3.1]{RuanHMinimal}である．
同論文は homothetic toric potentials と一様 $C^\infty$ quasi-isometry の下で，
potential の torus 平均，Hamiltonian-gradient flow，および陰関数定理を用い，
固定された縮小 polytope 上に H-minimal Lagrangian torus fibration を構成する．
Theorem~4.2 は simple toroidal canonical degeneration の最大退化点へ消える
minimal Lagrangian torus を与える．従って平均化と Moser 法の組合せ自体を
本稿の新規な方法とは扱わない．
本稿の定理\ref{ke:short-torus-averaging}は，$C^0$ の tensor 比較と短い作用軌道だけから，
変位と ambient fiber diameter を直接評価する弱い入力・弱い結論の定量化である．
H-minimal 性，special Lagrangian 性，変形後の intrinsic fiber metric は結論しない．

別研究\cite{KECubicalCurrent}の diameter-scale 候補は
\[
 R_t=\log(-\log|t|^2),\qquad
 (X_t,R_t^{-2}g_t)\longrightarrow(\mathcal C_d,d_\square).
\]
この収束は同稿の粗い比較仮定 $\mathrm{(CP)}_d$ と，
増大する積領域の一様比較・整合性を課す $\mathrm{(PBE)}_d$ の下でのみ主張されている．
さらに我々の fibration による収縮と解釈するには，適切に選びパラメータ化した
$\pi_t:X_t\to B_d(T)$ と系\ref{topo:KE-cubical-comparison}の $p_T$ について，
KE 側の近似写像 $\Phi_t$ が
\[
 \sup_{x\in X_t}d_\square\bigl(\Phi_t(x),p_T\pi_t(x)\bigr)\longrightarrow0
\]
となることを追加で示すのが一つの十分条件である．
命題\ref{topo:ke-map-compatibility}は，この条件から正則・特異 fibers の
\emph{ambient} 直径収縮と距離歪みの評価を導く初等的な比較原理である．
本稿はこの条件付き定理を引用して収束を結論する経路を採らず，
同じ $f_t=p_T\pi_t$ の全域上下界を直接証明する．
従って $\mathrm{(CP)}_d,\mathrm{(PBE)}_d$ 自体の証明を主張しない．
$d=5$ では $p_T$ は同相であり，$d\ge6$ で潰される内部木上の全 fiber 合併は
命題\ref{ke:collapsed-tree-full-fiber-collapse}で制御する．

三重点の積座標では，三枚の cubical squares を一つの toric hexagon によって
全閉領域上で同時に実現できる（命題\ref{ke:triple-hexagon}）．
定理\ref{ke:triple-moving-fibration}は，
$\ell_a=-\log|\zeta_a|^2>m_t/2$ の全 phase を含む領域で
actual KE metric と cusp Hermitian model の一様 tensor 比較を仮定し，
$\min\ell_a>m_t$ における Lagrangian fibration の
ambient fiber diameter $O(m_t^{-1})$ と，
raw logarithmic-depth map との誤差 $O((1+\log m_t)/R_t)$ を与える．
例えば $m_t=R_t$ と選べる．この tensor 比較は，体積密度の上下界だけからは従わない．
関連 KE 原稿\cite{KELocalCurrent}の
\texttt{coarse:prop:density-counterexample} は，滑らかな負 KE 計量に限っても
その一般的な推論が成立しないことを示す．
固定された affine simplex の compact subset は径の尺度では一つの triple vertex に縮むので，
既知の固定 compacta 上の局所収束だけで上の動く領域の仮定を埋めることもできない
（命題\ref{ke:fixed-simplex-scale}）．
この解析的比較は complete simplex model と境界体積比較によって検証できる．
系\ref{ke:actual-closed-corner}は任意の固定 $d\ge5$ で，三枚の閉 square 全体に一つの
actual Lagrangian torus bundle を構成し，$m_t=R_t$ のとき fiber 自身の intrinsic diameter も
$O(R_t^{-1})$ と評価する．
さらに定理\ref{ke:closed-corner-gh}は，同じ Moser flow と同じ fibers を用いて，
三枚の閉 square 自身の内在距離への projection distortion を
$O((\log R_t)/R_t)$ と評価する．
系\ref{ke:closed-corner-protected-gh}により，$\min\ell_i>R_t/2$ の保護領域内の
全経路を許してもこの評価は保たれる．
模型側の入力は，関連 KE 主稿の係数 $1$ の rank-one Schwarz 評価，
全正深度の二重 cusp 漸近，および clipped depth map による全経路の長さ下界である．
これに有限 square gallery の持ち上げを合わせ，model metric の Lie derivative 評価で
actual Lagrangian bundle へ移す．
標的の距離は閉 corner の内在距離として定義されるが，命題\ref{topo:coordinate-retraction}の非伸長 retraction により，全複体 $\mathcal C_d$ の距離の制限と等しい．
この局所距離評価だけから全曲面の shortcut の制御は得られない．

その相対接着に向け，命題\ref{ke:relative-fiber-preserving}は，
元の phase tori が既に actual Lagrangian である飽和 collar 上で
averaging の chain primitive による Moser flow が fiber 分割を保つことを示す．
Metric の torus 不変性は不要である．
Actual neck との orbit partition と quantitative buffer は後述の別の構成で検証し，
この命題だけから既存の大域 fibration への接着を結論しない．

同じ normal circle を持つ neck と triple collar を扱うため，
定理\ref{ke:staged-neck-averaging}では normal average $\eta=A_1\omega$ を先に作り，
その後に残る円を平均する．第一の map を $\phi$，第二を $\psi$ とすると，
$\psi$ は normal circle と可換であるから，$\Theta=\phi\psi$ の各 torus は
neck と共通の actual action $\phi\rho_1\phi^{-1}$ を含む．
単一の primitive $\omega=d\lambda$ を固定すれば
$\Theta^*\lambda=A_2A_1\lambda+dq$ であり，全整数 cycle の作用周期が一致する．
命題\ref{ke:staged-hamiltonian-comparison}は，平均順序や整数基底の違いを，
共通の buffer 上の Hamiltonian isotopy によって比較する．
Exact symplectomorphism であることだけを Hamiltonian 性の根拠にはせず，
primitive を exact 差まで保つ path homotopy を明示する．
系\ref{ke:staged-closed-corner-gh}は，この構成にも
$O((\log R_t)/R_t)$ の局所 GH 誤差を与える．
従来の一段階 bundle と同一の fiber 分割を主張するものではない．

Source-free な固定 compact base 上では，関連 KE 主稿\cite{KELocalCurrent}の
\texttt{pair:prop:bundle-conic} が射影的 normal-bundle coordinates から一つの正則円作用を与え，
\texttt{pair:cor:global-primitive} と \texttt{pair:cor:actual-common-moment} が
global primitive と円平均後の共通 moment interval を与える．
後者は固定 end での KE 収束と cusp 漸近を使う．
本稿の定理\ref{ke:fixed-base-neck-moser}は，capped density の大域下界と
Chern--Lu 評価を組み合わせ，固定残余基底上の moving normal strip で
actual 二側比較・短円・共通 Moser source と像を与える．
命題\ref{ke:fixed-base-moment-connectors}では同一 normal moment 上の残余点を
有界長で接続し，その level 全体の正規化 ambient 径 $O(R_t^{-1})$ を得る．
固定残余基底では source を含む一つの作用で actual neck と閉じた pinched torus も
構成する（定理\ref{ke:projected-source-neck}）．
補題\ref{ke:source-reduction-pinched-fiber}では，固定点での smooth exact transport と
canonical moment の零への規格化を確認し，位相的な零 reduction と ambient の
滑らかな Lagrangian strata を区別する．Normal orbit は primitive な消滅cycleであり，
閉じた source fibre は focus--focus 点を一つ持つ pinched torus となる．
Moment に依存する残余関数 $A(\mu,p)$ を用いた写像の非退化条件は
$d_pA(0,q)\ne0$ である（系\ref{ke:source-focus-transversality}）．
この写像の条件と葉の分割の型を分け，後続の相対補間では source の submersion germ を保つ．
Focus の局所型と monodromy 自体は既知の結果\cite[\S\S2--4]{ZungFocus}であり，
ここでは actual KE neck の共通作用・buffer・全閉 fibre へ適用する条件を実証する．

三重点への接続には，節\ref{ke:neck-corner-transition-section}で
固定比に分離した三つの arm を用いる．射影 Morse 座標と一つの共通 triple 座標の
対数比は，全 phase・全固定階の cusp 微分に関して指数的に小さい．
積を保つ smooth 補間は非正則になり得るが，平均形式と凸 path の tame 性を保つ．
Canonical primitive の base-only closed 差と exact 差を分離することで，
normal period の厳密な一致も得る．第二 period は別の条件である．
さらに節\ref{ke:relative-corner-section}で比較 Hamiltonian 自体を評価し，
三つの端に指定した normal-first 平均を一つの actual corner bundle へ相対統合する．
定理\ref{ke:slow-cut-closed-corner-gh}は
$m_t\to\infty$，$\log m_t=o(R_t)$ の遅い cut に対しても，
同じ corner 射影の局所 GH 近似を保証する．
節\ref{ke:neck-corner-exhaustion-section}では，固定残余 compact，normal cut，
共通 moment 幅，退化 parameter の閾値の順に選ぶ．Centered normal flow の自然性と
moment-dependent reduced foliation により，moving end の全軌道を共通領域に収める．
系\ref{ke:neck-corner-saturated-union}は，closed corner と切詰めた neck の
全ての実際の交わりで全 fiber が一致することを示し，それらの和集合上の
actual KE--Lagrangian 部分 fibration を与える．
補題\ref{ke:fixed-marked-tree-exhaustion}は，source の小円板での一価な isotopy と，
固定 compact の外に支持を持つ end 補正を用い，全 stage と全ての小さい複素 $t$ で
同じ marked pants 木系を選ぶ．これは全 torus cycle の parameter monodromy の
自明化や，異なる $t$ の nodal neighborhoods の semi-global symplectic 同値を意味しない．
後述の定理\ref{ke:actual-global-fibration}は，これを component cores へ相対接着して
全曲面を覆う．その同じ fibration の全曲面距離比較は，
後述の global clipped-depth map と共通 moment の接続で証明する．

Source の計量的な問題については，節\ref{ke:source-connectors-section}で
固定中央 rescaling の volume density を $L_t^{-2}$ と両側から比較する．
Canonical potential は一様に有界で，中央 moment は $O(L_t^{-1})$ にある．
Ricci 下界と Laplacian の障壁により actual inradius $O(L_t^{-1/2})$ を得て，
中央点を oriented arm の断面へ有界長で接続する．
系\ref{ke:fixed-base-central-band}は source-free な中央 moment 帯の全点を，
命題\ref{ke:source-central-band-reduction}は source 近傍の中央帯を固定 core へ，
それぞれ有界長で結ぶ．さらに固定 source tube 上の正の canonical potential の
局所対数勾配から，中央点の actual 球が normal 深度 $\kappa L_t$ 以上に留まることを示す．
距離函数の complex Sobolev norm に Dinh--Marinescu--Vu の
Moser--Trudinger 不等式\cite[Theorem~1.2]{DMVMT}を適用し，
logarithmic cover の周期数 $\asymp L_t$ を用いると，この広い中央帯で
actual 体積の Gaussian tail を得る．Bishop--Gromov 比較と合わせて，
定理\ref{ke:source-central-diameter}は固定点を含む中央 source の全点径を
$O(1)$ に抑える．固定残余基底では中央 moment 帯と閉じた pinched torus
全体にも適用できる（系\ref{ke:fixed-base-source-central-band}）．
用いた Moser--Trudinger 不等式自体は既存結果であり，ここで検証するのは
退化する実計量の距離函数への一様な適用，周期因子と actual 球の領域条件である．
積分測度には固定した補助 Euclidean potential を用い，actual potential の
一様 Hölder 連続性は仮定しない．
Potential の一様連続度による局所径評価\cite[Proposition~1.4]{GGZDiameter}は
同じ $O(1)$ の径を与える別の十分条件だが，その連続度はまだ得ていない．
本稿の全点径評価はその modulus を使わない．

Source の oriented arm では，normal covector の Schwarz 不等式に
固定 normal 境界・source-free な残余境界・moving normal 境界の上界を与える．
最後の境界では正の potential により orientation を保つ固定半径の actual 球を取り，
局所 Schwarz 評価を用いる．最大値原理で normal covector を全 arm へ制御し，
射影下界と体積上界を合わせて source 軸を含む二側 tensor 比較を得る
（命題\ref{ke:source-arm-uniform-tensor}）．
固定 normal 座標での残余水平経路と，最後の normal 深度補正により，
全共通 moment で同じ moment の残余境界へ $O(1)$ で退出できる
（定理\ref{ke:source-all-moment-exit}）．
中央点には短い gate 退出を先に使うので，この直接証明は上記 Gaussian 評価にも
依存しない（補題\ref{ke:source-direct-central-exit}）．
Gaussian 体積評価は独立な補助結果として保持する．
固定残余基底における同一 moment の全 fibers の合併も，
系\ref{ke:source-inclusive-same-moment-collapse}により収縮する．

成分 core については節\ref{ke:component-core-transfer-section}で，
fixed-gauge smooth KE 収束から canonical primitive の差を明示し，
固定 compact から actual KE 形式への小さい symplectic 移送を構成する．
境界の点を固定する相対移送には $H^1$ の制限写像で記述される flux 条件があり，
normal fiber periods だけではこれを検出できない．
一般位置の $m=d-1$ 直線補空間の compact core $K$ と $S=\partial K$ については，
Cohen--Suciu\cite[Proposition~2.5 and \S6.9]{CohenSuciuBoundary}の既知の境界 cohomology を用い，
\[
 H^1(S;\RR)/\operatorname{im}H^1(K;\RR)
 \cong H^2(K,S;\RR),\qquad
 \dim H^2(K,S;\RR)=\binom{d-2}{2}
\]
と計算できる．命題\ref{ke:component-flux-space}はこれを有限個の boundary periods による
相対 Moser の判定へ適用する．系\ref{ke:nonextendible-small-core-germ}では，
同じ actual KE 二形式を保つ小さい collar germs でも各非零障害方向を実現する．
平均形式も parameter 方向に Moser 移送することで exact collar germ も得る．
一方，定理\ref{ke:actual-global-fibration}の構成では，actual link の葉を先に固定し，
同じ actual 二形式の primitive を相対不変化して内部へ延ばすため，
点固定の相対 germ を追加入力とする必要はない．


開いた二次元セルの固定 compact 部分では，さらに明示的な product model と局所 GH 極限を得る．
関連 KE 原稿\cite{KELocalCurrent} の \texttt{simplex:cor:actual-dominant-product} は，
三重点の独立深度 $a,b\in[\tfrac12L_t^\eta,2L_t^{1-\eta}]$，$0<\eta<1/2$ において，
actual metric と product-cusp metric の全固定階の差を $O(L_t^{-\eta})$ と評価する．
証明では complete simplex と product domains の包含によって model volume ratio を $1$ に近づけ，
全深部三重点の actual semi-flat 比較と合わせ，
Savin\cite[Theorem~1.3]{SavinSmall} の小摂動正則性から tensor 比較を得る．
本稿の定理\ref{ke:open-cell-collapse}はこれを用い，
$[\eta,1-\eta]^2$ 上の actual KE--Lagrangian torus fibers の intrinsic diameter を
$O(L_t^{-\eta})$ とし，Moser 像領域の内在距離について平坦長方形への GH 収束を証明する．
Projection distortion も $O(L_t^{-\eta})$ である．
全曲面内の ambient distance は領域外の近道を許すため，この局所定理だけから
大域 GH 収束や既存 fibration の全 fibers に関する収縮を結論しない．

\paragraph{固定成分 core と target の距離．}
節\ref{ke:component-core-geometry-section}では complete KE 形式に対し，
任意の固定 compact 上で形式と primitive を保つ smooth 正則化と，
正 residue の有限型 Weinstein core を構成する．既存の相対 SNC 定理を用い，
その小さい飽和近傍を Lagrangian に filling する．

全域への相対構成は節\ref{ke:component-link-cover-section}--
\ref{ke:actual-core-gluing-section}で行う．Period $1$ の正の meridian に対し，
$M_a+2\pi w_a$ が既存の向き付き canonical normal moment と厳密に一致することを，
volume-frame ratio の整数 winding と exact averaging transport から示す．
$N_j\iota_j\gg1$ の深さで normal moment faces を取り，corner では二つの
moment cuts を smooth monotone minimum で丸める．
固定段階でも，さらに深い tube 上で同じ centered flow の逆像 buffer を使い，
fixed-shell moment の上界を加えると，既存の部分 fibration が中央の全特異 locus の
固定近傍を覆うことが分かる．その補集合は各 component の固定 compact に入る．

Rochon--Zhang の全 cusp 微分評価\cite[\S5, (5.3)--(5.6)]{RochonZhang}と
三重点の全固定階比較から，actual generators の $C^1$ smallness を導く．
Actual symplectic circle に対する centered scalar average は primitive を exact に補正し，
common normal action に不変な cutoffs で貼り合わせても
$Z^{\rm ad}\Psi>0$ を保つ．これにより既存 core を変えず，closed link 全体に外向きな
Liouville field を得る．Saddle を動かす相対葉族補間は，
各 critical pants と全 regular annuli に一つの vector field を作ることで証明する．
この補間は非零 moment band に置くので source の全 pinched tori は保たれる．

境界から内部への cone extension 自体は Evans--Mauri の既存結果
\cite[Lemma~2.16]{EM}に基づく．本稿では contact form を替えずに compact domain 上で用い，
actual link の既存葉，normal/corner overlaps，source を避けた相対補間，
全 fiber 単位の切詰めを合わせて検証する．
従って全 $X_t$ を覆う actual KE--Lagrangian fibration が得られる．
内側の primitive は compactly supported exact correction により Morse--Smale に選び，
index 順の閉 filtration によって skeleton の弱い isotropic stratification を確保する
（補題\ref{ke:compact-relative-morse-smale}）．
境界の contact form と既存 normal/corner actions は変えない．
各 component の変更領域は深度 $CN_j$ 以下の固定 compact image にあり，
cusp の動径経路により接続費用を $C(1+\log(N_j+1))$ と評価できる
（命題\ref{ke:component-logarithmic-connectors}）．
従って十分遅い対角選択でその正規化 ambient 径を消せる．
この定理自体の結論は保持 corner の局所構成を伴う全域の存在であり，
距離歪みの消失には次の追加評価を用いる．
この相対構成全体の一般的新規性は未確認である．

Target の組合せ型は既知の cubical barycentric subdivision\cite{BabsonBilleraChan}であり，
その構成自体の新規性は主張しない．命題\ref{topo:coordinate-retraction}は，unit-square 長さ距離に対する座標部分複体の等長性を与え，
閉 corner の target 距離を全体距離へ同定する．
定理\ref{topo:diameter-six-labels}は有限幾何のみで，六次以上の target の
厳密直径 $2\sqrt5$ と全最大対を決定する．
これらは actual 曲面の全域 shortcut を制御する結果とは区別する．

\paragraph{全曲面での距離下界と同じ基底値．}
Source を含む固定 residual domain 上に正則 normal frame $x$ を取り，
$s=-\log|x|^2$，$U=|d\log s|_{g_t}^2$ とする．
Bochner 計算から $\Delta\log U\ge U-1$ が得られる．
これは Schwarz 型評価\cite{YauSchwarz}の rank-one 版であり，
本稿では実計量と複素 Laplacian の規格化を含めて係数を直接計算する．
固定外側 shell，残余基底の境界，moving central shell の三つの境界を検証し，
\[
 B=1+\varepsilon_j+A_j\mathfrak n_j/s+C_js/(\kappa L_t)
\]
を障壁として用いる．二つの cut を十分離すことで，source を除くことなく
active な $|d\log\ell_i|_{g_t}$ を $1+o(1)$ に抑える
（定理\ref{ke:sharp-global-one-active}）．
体積密度の二側比較だけから sharp tensor bound を推論してはいない．

Global clipped-depth map $\Theta_t$ は，高々三つの非零 coordinates と
少なくとも一つの plateau を持つ．二つの active coordinates では全深部 triple comparison，
一つの場合は上の scalar estimate を用いる．
得られる長さ下界は全曲面の任意の経路に直接成立し，領域の出入りの回数に依存しない．
その後で，保持 corner では旧 softmax と短い Moser 移送を比較する．
Neck の任意の点には，同じ canonical moment を持つ旧 corner 境界上の参照 torus がある．
Source と branch を含む同一 moment 集合の接続評価と，既に証明した
$\Theta_t$ の Lipschitz 評価により，その参照 torus での値比較を全 neck へ移す
（補題\ref{ke:same-moment-corner-reference}）．
Core の像の小ささと合わせて
$\sup_{X_t}d_\square(\Theta_t,p_T\pi_t)=o(1)$ を得る．
ここで用いる接続評価は局所解析から独立に証明してあり，全域の距離上界を先取りしない．
Clipped depth が個々の Lagrangian fiber 上で一定であるとも仮定しない．

\paragraph{基底の接着と全曲面での距離上界．}
旧 corner の境界 torus と新しい cellular branch locus は同じ strata ではない．
命題\ref{ke:global-base-parameterization}は保持 corner を少し内側へ写し，
その外周に neck strips と component sectors を配置する．
旧 softmax 値との差を $o(1)$ に保ち，内部木の全枝に一つの canonical moment を使う．
この同じ moment 上の source を含む接続定数 $C_j^{\rm neck}$ は，
component cut $N_j$ より先に固定できる．Neck の component 側の尾部も深度
$CN_j$ 以下に入ることを別に示す．従って
$C_j^{\rm neck}\le1+\log(N_j+1)$ となるよう cut を選べば，
全接続に一つの予算 $H_j=C(1+\log(N_j+1))$ を使える．
Target の最短路を高々 $60$ 本の square 線分に分け，共通 moment の portals で繋ぐと
同じ $f_t$ の全曲面距離上界を得る．その誤差は
$C_d\{b_j^{-1}+(1+\log R_t+\log(N_j+1))/R_t\}$ 以下である．
同一 moment の参照点を使った下界でも，残余接続費用をこの同じ cut に吸収できる．
式\eqref{ke:eq:global-realization-explicit-cut-budget}は，上下両側を合わせた distortion を
固定族の一つの係数と，$b_j,\varepsilon_j,\mathfrak m_j,K_j,N_j$ の明示式で抑える．
これは許容 cut の下での評価であり，残余定数の増大率や最適収束速度を主張しない．
定理\ref{ke:global-metric-realization}では有限個の stage 定数を先に選び，
共通の遅い対角化で上下両側の誤差を零にする．
これらの source 評価，scalar 障壁，基底の相対 parameter 化と全域距離比較が，
本稿で既存の局所理論に加えて証明する部分である．
全体の先行研究に対する新規性の確定と包括的な独立検証は，公表前の課題である．

\subsection{他の tropical 構造との比較と残る問題}

Cartwright の tropical complex は，退化の dual complex に intersection numbers を加えて，
線形関数の層を指定する構造である\cite[Definitions~2.5, 2.8; Construction~3.2]{CartwrightComplexes}．
作用周期から作る本稿の atlas と同一視することはできない．
実際，ordered small resolution の最初の component における local intersection matrix は，
全平面の場合 $J_{d-1}-dI$ で負定値となり，tropical complex の local positivity condition を満たさない
（命題~\ref{atl:cartwright-distinction}）．
同論文 Remark~3.4 も，Gross--Siebert 型の affine structure との特異集合の違いを区別している．
本稿でも，focus points を残した atlas に compact smooth tropical manifolds の大域定理を
追加の検証なしに適用しない．

また，標準 fan の chart と特定の algebraic filling との同定は別の条件である．
Kapranov\cite[Theorem~4.3.3]{KapranovChow}の blow-up 記述によれば，
$\overline M_{0,5}$ の boundary は十本の $(-1)$-curves の Petersen 配置である．
五直線 complement の模型は common-circle determinants と boundary homology によって，
この標準 $M_{0,5}$ filling と区別される
（命題~\ref{trop:divisor-recovery}，系~\ref{trop:not-standard-M05}）．
二つの boundary の $b_1$ はそれぞれ十と十一である．
従って作用適合 Petersen chart の存在から，その filling の置換を結論してはいけない．

全平面五次の高分母 chart の同時接着とその支持 graph の問題は，
companion 論文\cite{QuinticArithmetic}で，整数データの必要条件と
幾何学的な実現条件を区別して検討する．そこでは非空の高分母支持が葉を持たず，
最小次数三以上・高辺八本以上の三つの五頂点 graph 型，二十六標識に限られることを示す．
非空支持は三点連結であり，単位分母の辺は matching をなす．
特に一つの core に接する二本の branches の分母が一なら，全分母が一となる．一方，
残る整数接着の存在・一般の分母剛性は未解決である．
同 companion 論文では，全ての pair products が正である四ray整数データの正値性を，
完全 matching との同値性で判定する．
単位枝の node 配分は $\{5,0,0\}$ または $\{4,1,0\}$ に限られ，
二単位枝を持つ $G_{5,14}$ では後者だけが残る．
これらは局所の正面積と整数接着に関する必要条件であり，残る三型の大域実現を主張しない．
本稿の構成に続く幾何的な問題として，atlases の同値関係と変形空間，
および core を替える相対 surgery が残る．
四 puncture の内部 cones の格子 index を $m^-,m^+$ に緩める場合も，
$q_Aq_B\mid(m^-+m^+)$ という必要条件が残る．
一般次数における非 matroidal links を，より広い tropical models で扱う可能性も残る．
hyperplane class 以外の Kähler classes，special Lagrangian 条件，本稿以外の退化族における canonical metric の collapse は，
本稿の存在定理から独立した課題として扱う．
Gross--Siebert\cite{GrossSiebert}の affine geometry からの退化再構成との比較にも，
本稿で得た実際の作用周期に加えて，必要な追加データを明示する必要がある．
